\documentclass[20pt,a4paper,landscape]{extarticle}
\usepackage[margin=1.2cm,top=0.6cm,bottom=0.8cm,left=1.0cm,right=1.0cm]{geometry}
\usepackage{amsmath}
\usepackage{amssymb}
\usepackage{array}
\usepackage{enumitem}
\usepackage{multicol}
\usepackage{xcolor}
\usepackage[sfdefault,lf]{FiraSans}
\renewcommand{\familydefault}{\sfdefault}
\setlength{\parindent}{0pt}
\pagestyle{empty}
\everymath{\displaystyle}
\setlist[itemize]{leftmargin=1.3em,itemsep=0.35em,topsep=0.2em}
\setlist[enumerate]{leftmargin=1.5em,itemsep=0.45em,topsep=0.2em}
\setlength{\columnsep}{1.2cm}

\definecolor{mmblue}{HTML}{1F4E79}
\definecolor{mmsoft}{HTML}{EEF4F8}
\definecolor{mmwarn}{HTML}{A63D40}
\definecolor{mmlink}{HTML}{315C2B}

\newcommand{\NotesTitle}[1]{{\fontsize{28}{31}\selectfont\textbf{#1}}\par\vspace{0.25em}\color{mmblue}\rule{\linewidth}{1.2pt}\color{black}\vspace{0.55em}}
\newcommand{\SectionTitle}[1]{{\Large\textbf{\textcolor{mmblue}{#1}}}\par\vspace{0.25em}}
\newcommand{\KeyIdea}[1]{\colorbox{mmsoft}{\parbox{0.97\linewidth}{\vspace{0.35em}\textbf{Key idea:} #1\vspace{0.35em}}}\par\vspace{0.5em}}
\newcommand{\WorkedExample}[2]{\fbox{\parbox{0.97\linewidth}{\vspace{0.35em}\textbf{#1}\par\vspace{0.35em}#2\vspace{0.35em}}}\par\vspace{0.55em}}
\newcommand{\CommonMistake}[1]{\fbox{\parbox{0.97\linewidth}{\vspace{0.35em}\textbf{\textcolor{mmwarn}{Common mistake}}\par\vspace{0.25em}#1\vspace{0.35em}}}\par\vspace{0.5em}}
\newcommand{\TryThese}[1]{\fbox{\parbox{0.97\linewidth}{\vspace{0.35em}\textbf{\textcolor{mmlink}{Try these}}\par\vspace{0.25em}#1\vspace{0.35em}}}\par}
\newcommand{\StepLine}[2]{\textbf{#1.} #2\par\vspace{0.35em}}

\begin{document}
\sffamily
\boldmath

\NotesTitle{Topic 1 LO1 — Place Values and Decimals}

\begin{multicols}{2}

\KeyIdea{A digit's value depends on its place. The same digit can mean very different amounts in different positions.}

\SectionTitle{Steps}
\StepLine{1}{Look at the digit you care about.}
\StepLine{2}{Name its place: ones, tens, tenths, hundredths, and so on.}
\StepLine{3}{Write its value, not just the digit.}
\StepLine{4}{If comparing numbers, check from left to right.}

\columnbreak

\SectionTitle{Place-value reminder}
\begin{itemize}
  \item In \(347\), the \(4\) means \(40\).
  \item In \(5.3\), the \(3\) means \(0.3\).
  \item In \(2.74\), the \(7\) means \(0.7\).
  \item In \(2.74\), the \(4\) means \(0.04\).
\end{itemize}

\WorkedExample{Example 1: value of a digit}{
Find the value of the \(6\) in \(862\).
\[
862 = 800 + 60 + 2
\]
So the value of the \(6\) is \(60\).
}

\WorkedExample{Example 2: decimal place value}{
Find the value of the \(7\) in \(2.74\).
\[
2.74 = 2 + 0.7 + 0.04
\]
So the value of the \(7\) is \(0.7\).
}

\end{multicols}

\GreenDivider
\CommonMistake{Saying the value is just the digit. In \(862\), the digit is \(6\), but its value is \(60\). In \(2.74\), the digit is \(7\), but its value is \(0.7\).}

\TryThese{%
\begin{enumerate}
  \item What is the value of the \(4\) in \(347\)?
  \item What is the value of the \(5\) in \(5.3\)?
  \item What is the value of the \(4\) in \(3.48\)?
\end{enumerate}}

\newpage

\NotesTitle{Topic 1 LO1 — Place Values and Decimals}

\begin{multicols}{2}
\SectionTitle{Comparing numbers}
When two numbers have the same whole-number part, compare the decimal places next.

\WorkedExample{Example 3: compare decimals}{
Which is greater: \(0.8\) or \(0.08\)?
\[
0.8 = 8\text{ tenths}
\]
\[
0.08 = 8\text{ hundredths}
\]
A tenth is bigger than a hundredth, so
\[
0.8 > 0.08
\]
}

\WorkedExample{Example 4: order numbers}{
Order \(3.2,\ 2.9,\ 3.02\) from smallest to largest.
\[
2.9 < 3.02 < 3.2
\]
So the order is:
\[
2.9,\ 3.02,\ 3.2
\]
}

\columnbreak

\SectionTitle{Partitioning numbers}
Split a number into place-value parts.

\WorkedExample{Example 5: whole number}{
Partition \(407\).
\[
407 = 400 + 0 + 7
\]
Usually we write this as:
\[
407 = 400 + 7
\]
}

\WorkedExample{Example 6: decimal number}{
Partition \(3.48\).
\[
3.48 = 3 + 0.4 + 0.08
\]
That is 3 ones, 4 tenths, and 8 hundredths.
}
\end{multicols}

\GreenDivider
\CommonMistake{Thinking more digits means a bigger decimal. For example, \(0.08\) is not bigger than \(0.8\). It is smaller because 8 hundredths is less than 8 tenths.}

\TryThese{%
\begin{enumerate}
  \item Which is greater: \(45\) or \(54\)?
  \item Order from smallest to largest: \(0.5,\ 0.05,\ 0.15\)
  \item Partition \(4.7\)
\end{enumerate}}

\end{document}
